\section{Banked constructions and attainable values}
\label{val:section}

A local separation identity must remain valid when the positions on which it
acts are correlated with positions to be processed later.  It is also necessary
to distinguish a distribution in the coordinates of a tensor from the
aggregate distribution after physically permuting its sides.  We make both
requirements explicit.  This gives the finite objects whose attainable values
will enter the differential comparison.

\subsection{Canonical banks and preservation of unresolved tensors}
\label{val:banks}

Let $\mathcal S=\{200,020,002,011,101,110\}$ and
$\mathcal W=\{100,010,001\}$.  The latter is the support of the ordinary
three-state tensor
\[
 \mathsf W=x_1y_0z_0+x_0y_1z_0+x_0y_0z_1.
\]
An element of either alphabet is a triple of local grades.
We use $\mathcal B$ to denote either alphabet in definitions
that apply to both.  Write $\mathfrak S_3$ for the group of permutations of the
three tensor sides.  A \emph{banked input of size $N$} has $N$ occurrences in
each orientation $\sigma\in\mathfrak S_3$, for a total of $n=6N$
occurrences.  Its original component words are
\[
 w=(w_{\sigma,t})_{\sigma\in\mathfrak S_3,\ 1\le t\le N},
 \qquad w_{\sigma,t}\in\mathcal B.
\]
The triple $w_{\sigma,t}$ is expressed in canonical coordinates: its physical
grade triple is $\sigma w_{\sigma,t}$.  Thus physical side $j$ observes the
canonical coordinate $\sigma^{-1}(j)$.  Both $T$ and the ordinary $\mathsf W$ tensor
are symmetric under these physical permutations, so this banking does not
change the underlying input tensor.  Its purpose is to record types and
permitted transposes accurately.

A banked word has exact law $\alpha\in\Delta(\mathcal B)$ if, separately for
each $\sigma$, every symbol $s$ occurs $N\alpha_s$ times.  In particular,
$\alpha$ is not the histogram obtained by forgetting the bank orientations.
All normalizations below divide by $6N$.

We allow additional, unresolved local indices.  To state the preservation
requirement precisely, for every original component word $w$ attach an
arbitrary three-tensor $P_w$ on these additional indices.  The tensors $P_w$
may share their variable spaces, may correlate all the positions, and may be
zero.  Denote the resulting input by $\mathcal T(P)$: its restriction to the
original word $w$ is the original scalar component at $w$ tensored with
$P_w$.  This description includes preceding local count restrictions by
setting the corresponding fibers to zero.  It does not assert that all such
families arise as independent products.

\begin{definition}[Strong banked realization]
\label{def:strong}
A strong realization consists of a finite auxiliary tensor $\mathcal A$, a
finite degeneration by three local maps, a set $\Lambda$ of direct labels,
and positive integers $a,b,c$, with the following properties.
\begin{enumerate}[label=\textup{(\roman*)}]
\item Every $\lambda\in\Lambda$ designates an original word $w(\lambda)$,
      and $\lambda\mapsto w(\lambda)$ is injective.  All these words have
      the same exact canonical law in each orientation bank.
\item For every family of unresolved tensors $P=(P_w)$ the same maps give
\begin{equation}
\mathcal T(P)\otimes\mathcal A
 \unrhd
 \bigoplus_{\lambda\in\Lambda}
 \bigl(P_{w(\lambda)}\otimes\langle a,b,c\rangle\bigr).
\label{val:strong-identity}
\end{equation}
      In this display the original local coordinates of $w(\lambda)$ are
      retained as well; only their fixed basis factors are suppressed.
      Original coefficients and unresolved local indices are unchanged.
\item The auxiliary input, its charged border-rank bound, and the output
      dimensions $a,b,c$ are independent of $\lambda$.
\end{enumerate}
The identity is required also for zero $P_w$.  Such an output is a formal
zero fiber and contributes no nonzero direct summand.  A realization used in
an attainable-value calculation must separately establish that all its
counted fibers are nonzero on its specified input.
\end{definition}

The maps in this paper retain each original local index and act on auxiliary
coordinates conditionally on that index and on already direct labels.
Consequently they commute with local projections depending only on the
original indices.  The same statement holds when the maps are expressed as
leading monomial degenerations.  This observation concerns locality; it does
not prove nonvanishing after an earlier projection.  Nonvanishing will be
checked by support identities for the algebraic primitives and by the
universal path condition in Lemma~\ref{lem:compile} for arithmetic-progression
(AP) count constructions.

Lemma~\ref{lem:group} supplies \eqref{val:strong-identity} term by term.
Indeed its sphere and slice argument decides the retained auxiliary terms
without referring to any coefficient of $P_w$.  Local singleton projections
and direct matchings have the same preservation property.  The conditional
composition rules of Section~\ref{val:recipes} use only these identities;
they never presume independence of unresolved fibers.

For the banking operation of Lemma~\ref{alg:orientation-banks}, group the
original input occurrences of a finite schedule by their relative physical
orientation $\tau$.  The lemma requires their total canonical histogram
$H_\tau$ to be fixed independently of the direct history and shared by all
six globally permuted replicas.  Exact types of complete original words
supply these grouped histograms in the completion and parity recipes.
When an already strong child is physically permuted as a whole, each of
its original input banks supplies a fixed histogram by
Definition~\ref{def:strong}.  The letter at an individual position may
still vary between direct outputs.

In outer replica $\sigma$, the group of relative orientation $\tau$ has
absolute orientation $\sigma\tau$.  Each absolute bank therefore receives
one copy of $H_\tau$ for every $\tau$, and has histogram
$\sum_\tau H_\tau$.  Equal conditional pool sizes alone would not imply
this histogram identity.  The pool sizes instead fix how many child calls
and auxiliary supplies are used; the grouped types or strong child banks
supply the required letter counts.  A physical permutation of an
intermediate macro acts on every original occurrence inside it.

\subsection{Finite recipes and the closure of their values}
\label{val:recipes}

A \emph{structural recipe} is a finite expression describing local
projections and leading degenerations, direct-label-controlled maps, locally
imposed exact joint and conditional types, auxiliary GROUP calls, and tensor
products on specified collections of positions.  Its bank sizes, rational types, and
integer group parameters are finite.  A conditional choice is permitted only
after its controlling label is direct on all three sides.  Each byproduct
matrix multiplication factor is kept as a separate factor: later decisions
use original-word labels and do not read, identify, or spend its indices.

We use the following two classes of such recipes.
\begin{enumerate}[label=\textup{(\alph*)}]
\item The $\mathsf W$ class starts with singleton projections, binary deletion
      followed by GROUP, and the finite COMPLETE and PARITY compositions
      of Proposition~\ref{alg:recipe-composition}.  It allows rational
      time-sharing on separate positions, physical permutation banking,
      and finite AP count compositions of previously constructed recipes.
\item The six-state class starts with singleton projections, the direct
      matching construction of Proposition~\ref{prop:matching}, and the
      five-state face substitutions of Proposition~\ref{prop:face} using
      finite $\mathsf W$ recipes.  It has the same time-sharing, banking, and finite
      AP closure.
\end{enumerate}
Both classes also allow two explicit finite rearrangements.  First, one
may reindex the orientation banks by $\sigma\mapsto\sigma\tau$ for a fixed
$\tau\in\mathfrak S_3$, with the corresponding canonical letter relabeling;
Lemma~\ref{val:canonical-symmetry} proves its precise effect.  Second, one
may \emph{conditionally recompile} an earlier finite recipe: expand its
finite direct-label refinement tree, take finitely many replicas, impose a
rational exact type of complete terminal transcripts through its induced
conditional types, and replace corresponding GROUP calls by GROUP on the
longer conditional pools.  All other original support predicates are retained
on their old macroblocks.  This is a finite construction rule using the same
local primitives, not an operation on scores.  Its instances must preserve
Definition~\ref{def:strong} and nonzero support like every other recipe;
Lemma~\ref{val:recompile} proves those requirements for the uniform terminal
type used below.  Neither rearrangement authorizes an infinite composition.

Here an AP composition is the finite count-window construction of
Lemma~\ref{lem:compile}.  Its children are finite recipes, not calls to an
attainable-value function.  The full expression can therefore be expanded to
a finite tree.  Equivalently, take the union of all finite levels of this
construction rule.  The $\mathsf W$ class is defined independently of the six-state
class.

A recipe belongs to these classes only when its local identities establish
Definition~\ref{def:strong}, its counted original words survive all of its
restrictions, and its auxiliary profile and output dimensions are fixed
across its direct histories.  These are algebraic conditions, not inequalities
for numerical scores.  The primitive support proofs, exact conditional-type
composition, and Lemma~\ref{lem:compile} establish these conditions for the
listed generators.  In particular, introducing the AP closure here does not
assume the differential inequality that it will later prove.

The auxiliary tensor in every recipe is a product of finite group tensors.
We charge the product of their ranks, which is also their tensor product's
rank by Fourier diagonalization and flattening.  It is harmless to charge a
larger certified bound.  No tensor already present in an unresolved fiber is
included in the byproduct volume.

For a realization with $n=6N$ original occurrences, $L=|\Lambda|$ counted
nonzero direct summands, auxiliary rank bound $R$, and common additional
volume $v=abc$, put
\begin{equation}
 D=\frac{\log L}{n},\qquad
 A=\frac{\log R}{n},\qquad
 M=\frac{\log v}{n}.
\label{val:rates}
\end{equation}
For $0<\beta\le3$ and $\gamma\ge1$, respectively, the two scores are
\begin{equation}
 f_\beta=D-A+\frac\beta3M,
 \qquad
 u_\gamma=\gamma(D-A)+M.
\label{val:scores}
\end{equation}
These are formal tensor rates.  Their definitions impose no assumption on
$\omega$.  In particular, on the same realization
$u_{3/\beta}=(3/\beta)f_\beta$.

\begin{definition}[Attainable-value functions]
\label{val:def-value}
For each finite realization let $\alpha$ be its canonical bank law.  Take all
pairs $(\alpha,t)$ with $t$ no greater than its score, and take their ordinary
Euclidean closure in $\Delta(\mathcal B)\times\R$.  The upper boundaries of
these closed sets are denoted by $F_\beta(\alpha)$ for the six-state class
and $U^W_\gamma(q)$ for the $\mathsf W$ class.
\end{definition}

Allowing smaller scores makes these sets downward closed.  The closure allows
irrational laws and limits of finite recipes of increasing size; it does not
introduce an infinite composition acting on a finite tensor.  In particular,
if $F_\beta(\alpha)>c$, then some actual finite realization has score greater
than $c$, possibly at a law arbitrarily close to $\alpha$.

\begin{proposition}[Bounds, concavity, and continuity]
\label{prop:values}
For $0<\beta\le3$ and $\gamma\ge1$,
\begin{equation}
 0\le F_\beta(\alpha)\le H(\alpha)\le\log6,
 \qquad
 0\le U^W_\gamma(q)\le\gamma H(q)\le\gamma\log3.
\label{val:entropy-cap}
\end{equation}
Both functions are concave and continuous on their closed probability
simplices.  They have supporting affine functionals at every relative
interior point.
\end{proposition}

\begin{proof}
Every additional pairing has size at most the rank charged for its group
tensor.  Because byproducts remain separate, their volumes multiply, and
$M\le A$.  Source-word injectivity and exact types give
\[
 L\le\left(\frac{N!}{\prod_s(N\alpha_s)!}\right)^6
 \le \exp(6NH(\alpha)),
 \qquad D\le H(\alpha).
\]
The same formula holds with the $\mathsf W$ alphabet.  Thus $f_\beta\le D$ and
$u_\gamma\le\gamma D$, proving the upper bounds also after closure.

A specified original word can be retained by singleton projections.  Taking
one such word of any prescribed rational type in each bank gives $L=R=v=1$
and score zero.  Density supplies the lower bound zero at every law.

To mix two finite constructions, repeat them enough times that a rational
fraction $\lambda$ of the positions in each orientation is allocated to the
first and the remaining fraction to the second.  Tensor their maps, auxiliary
inputs, and outputs.  The new law and score are
\[
 \lambda\alpha^{(1)}+(1-\lambda)\alpha^{(2)},\qquad
 \lambda f^{(1)}+(1-\lambda)f^{(2)}.
\]
Their direct transcripts inject into concatenated original words, all bank
histograms and dimensions remain uniform, and all support conditions are
preserved.  Rational approximation and closure give concavity for real
$\lambda$ and limiting laws.

For upper semicontinuity, suppose $\alpha_k\to\alpha$.  The upper bounds in
\eqref{val:entropy-cap} give a bounded subsequence of values approaching the
limsup.  Approximate each value by a pair in the defining closed set.  Closedness
then places the limiting pair in that set, so its value is at most the value
at $\alpha$.  This also shows that each finite upper boundary is attained in
the closed set.

For lower semicontinuity at any $\alpha$, including at a face, set
\[
 \varepsilon_k=
 \max_{s:\alpha_s>0}
       \left(1-\frac{\alpha_{k,s}}{\alpha_s}\right)_+.
\]
Then $\varepsilon_k\to0$.  If $\varepsilon_k>0$, the vector
\[
 r_k=\frac{\alpha_k-(1-\varepsilon_k)\alpha}{\varepsilon_k}
\]
is a probability vector.  Indeed its coordinates are nonnegative, including
those where $\alpha_s=0$, and its total is one.  Concavity and nonnegativity
give $F_\beta(\alpha_k)\ge(1-\varepsilon_k)F_\beta(\alpha)$.
If $\varepsilon_k=0$, conservation of total mass implies
$\alpha_k=\alpha$.  This proves lower semicontinuity; the proof for $U^W_\gamma$
is identical.  A finite continuous concave function has a supporting affine
functional at each relative interior point, by separation of its hypograph.
\end{proof}

\begin{lemma}[Symmetry of canonical laws]
\label{val:canonical-symmetry}
For $\tau\in\mathfrak S_3$, define
$(\tau_*\alpha)(s)=\alpha(\tau^{-1}s)$.  Then
\[
 F_\beta(\tau_*\alpha)=F_\beta(\alpha),\qquad
 U^W_\gamma(\tau_*q)=U^W_\gamma(q).
\]
\end{lemma}

\begin{proof}
Start with a strong realization and move old bank $\sigma\tau$ to new bank
$\sigma$, reinterpreting its canonical symbol $s$ as $\tau s$.  The physical
grade is unchanged, because $\sigma(\tau s)=(\sigma\tau)s$.  This is a
fixed permutation of original occurrence positions, performed identically
on all three sides, together with a change of canonical bookkeeping.
Conjugate the old local maps by this position permutation.  Their
coordinatewise identities continue to hold for arbitrary unresolved tensors,
with the same relabeling of their original positions.  Every new bank has
canonical law $\tau_*\alpha$.  Source-word injectivity, nonzero support,
the auxiliary tensor, and all byproduct dimensions are unchanged.  Thus the
same $D,A,M$ are attained at the permuted law.  The inverse rearrangement
uses $\tau^{-1}$; taking closures proves both equalities.  This canonical
law symmetry is distinct from replicating a schedule in all six physical
orientations, which preserves its canonical law.
\end{proof}

\begin{remark}[A useful exact-law correction]
\label{val:dilution}
If a construction or a convex mixture has law $\widehat\alpha$, and a desired
law $\alpha$ has positive coordinates, put
\[
 \lambda=\max_{s:\widehat\alpha_s>0}
 \left(1-\frac{\alpha_s}{\widehat\alpha_s}\right)_+.
\]
For $\lambda>0$, the vector
$r=(\alpha-(1-\lambda)\widehat\alpha)/\lambda$ is a probability law.
Mixing the construction with singleton constructions of average law $r$
therefore gives $\alpha$ exactly, at loss at most $K\lambda$ if the old score
is at most $K$, where $K\ge0$.  Moreover
\[
 \lambda\le
 \max_s\frac{(\widehat\alpha_s-\alpha_s)_+}{\alpha_s}.
\]
When the data are rational, this is a finite rational time-sharing operation.
For limiting data it is justified by the defining closure.  This correction
will be used when converting approximate finite laws to exact target laws.
\end{remark}

\subsection{Conditional pooling and monotonicity in the formal parameter}
\label{val:pooling}

A finite GROUP call pays a small excess of auxiliary rank over its number of
labels.  Tensor-repeating that particular finite map does not remove its
normalized excess.  The following lemma replaces each GROUP call by a new
one acting on a long conditional pool.  It is the requisite justification for
monotonicity in $\gamma$.  The face substitution below evaluates the W value
at $\gamma=3/\beta$; monotonicity will let the certificate use a lower
estimate obtained at a smaller formal parameter.

\begin{lemma}[Recompiling a fixed finite recipe]
\label{val:recompile}
Fix a finite recipe on $n$ original occurrences with canonical law $q$, $L$
terminal labels, charged auxiliary rank $R$, and additional volume $v$.
Fix $\gamma\ge1$.  There is a sequence of recipes on larger exact banks,
with the same canonical law $q$, such that
\begin{equation}
 \liminf u_\gamma^{\mathrm{new}}
 \ge \frac{\gamma\log L-\gamma\log R+\log v}{n},
 \qquad
 \liminf(D^{\mathrm{new}}-A^{\mathrm{new}})\ge0.
\label{val:recompile-goal}
\end{equation}
All source-word restrictions of the old recipe hold separately in every old
macroblock of the new recipes.
\end{lemma}

\begin{proof}
Expand the fixed recipe into its finite sequence of conditional refinements,
including the finite recipes used at any earlier AP nodes.  Refine a history
whenever its next local operation, its informed pair, or its auxiliary
parameters differ.  Terminal histories can be padded by deterministic steps
so that every transcript is a leaf of one finite decision tree.  Every
intermediate label is determined by the designated original word: GROUP
labels encode its source label sequence, matching labels identify the word,
coarse flags are functions of it, and AP labels are fixed by its disjoint
count windows.  Thus this expansion adds no terminal multiplicity; its leaves
are exactly the $L$ original strong labels.  A GROUP
context refers to a particular such call with its preceding direct history
fixed.  If several pools are processed, order them and include their labels
in the history.  A physical orientation is part of the context when needed.

Give the $L$ complete terminal transcripts their uniform probability law.
At context $c$, let $w_c$ be its visitation probability, let $r_c$ be the
conditional law of its GROUP label, and write
\[
\begin{aligned}
 h_c&=H(r_c),\\
 K_c&=\text{number of available labels},\\
 R_c&=\text{charged group rank},\\
 m_c&=\text{old pairing size}.
\end{aligned}
\]
Repeated calls are regarded as distinct contexts.  All probabilities are
rational.  Every actually used label is among the $K_c$ available tags, so
\begin{equation}
 h_c\le\log K_c\le\log R_c,
 \qquad \log m_c\le\log R_c.
\label{val:group-inequalities}
\end{equation}
The last two inequalities follow because there are at most $|G|$ tags and at
most $|G|$ retained pairing coordinates in a group tensor of rank $|G|$.

Choose $B$ through sufficiently divisible multiples.  Across $B$ replicas of
the old complete recipe, retain precisely the exact type with $B/L$
occurrences of each terminal transcript.  This condition is implemented in
prefix order: at a prefix of probability $w$, impose the induced exact
conditional type of its next label on its $Bw$ replicas.  Each next type is
readable on the sides that already read that label.  For a direct matching
label or an AP target label, all three sides read it; for a GROUP label, its
two informed sides read it before the GROUP call.  Previously direct labels
identify the relevant pool on every side.  The conditional type restrictions
are therefore local.

At GROUP context $c$ replace the old $Bw_c$ copies of its auxiliary call by
one application of Lemma~\ref{lem:group} to the entire conditional label
sequence of that pool.  Its number of possible labels is
\[
 K'_c=\frac{(Bw_c)!}
           {\prod_s(Bw_cr_c(s))!},
 \qquad
 \log K'_c=Bw_ch_c+O(\log B).
\]
Choose the pairing size to be $K'_c$ and choose the efficient group parameters
of Lemma~\ref{lem:group}.  Consequently
\begin{equation}
 \log R'_c=Bw_ch_c+o(B),\qquad
 \log m'_c=Bw_ch_c+o(B).
\label{val:new-group-rates}
\end{equation}
A deterministic pool requires no auxiliary.  The number of contexts is
fixed before $B$ increases, so all their errors sum to $o(B)$.

Here are the support and resource details of this replacement.  The current
pool is selected by previously direct histories.  Both informed sides can
read its conditional label word independently, so the new GROUP call is
legitimate.  Its payload comprises every remaining original coordinate and
every old support predicate.  The pointwise identity preserves that payload.
The new direct tag identifies the entire conditional word on all three sides,
so each old per-macro GROUP label is locally recoverable from it.  Every
subsequent old label-controlled map can therefore be implemented using the
new tag; the old pairing register is replaced and is never consulted.
Other direct operations are retained on the old macroblocks, followed by the
specified conditional-type projections.  Thus each final tuple consists of
old admissible transcripts, and every old per-block count window,
COMPLETE/PARITY predicate, and singleton or matching restriction still holds.
All its original scalar fibers are nonzero.  The exact conditional counts
fix every pool and every auxiliary multiplicity independently of the new
output history.  The auxiliaries can be appended once and the maps used
block-diagonally on direct histories.  New pairing dimensions may differ
from the old ones; this does not affect a later map, since byproduct indices
are never used for a later decision.

The product of the conditional multinomial counts is the full multinomial
count of terminal transcript tuples.  Hence the new direct multiplicity is
\begin{equation}
 L'=\frac{B!}{((B/L)!)^L},\qquad
 \log L'=B\log L+O(\log B).
\label{val:new-label-count}
\end{equation}
No terminal transcript is lost beyond these prescribed types: non-GROUP
maps act as before, and each replacement GROUP retains every conditional
word of its selected type.  Conversely, these types specify exactly the
chosen leaf counts.  This also proves injectivity into original words.
Each old terminal word has the same law in every canonical bank.
Concatenate the $B$ already banked replicas within corresponding original
orientation banks, retaining every old orientation label and internal
relative permutation.  No additional outer replicas are introduced at this
stage.  Each resulting canonical bank therefore has the old law $q$.
The construction respects Definition~\ref{def:strong}, not just a total
histogram after forgetting orientations, with exactly $Bn$ original
occurrences as used in the displayed counts.

The chain rule for the uniform terminal transcript gives
\begin{equation}
 \sum_{c\text{ a GROUP context}}w_ch_c\le\log L.
\label{val:group-chain-rule}
\end{equation}
The omitted contributions are conditional entropies of other direct labels
and are nonnegative.  Equations~\eqref{val:new-group-rates} and
\eqref{val:new-label-count} now prove the second assertion in
\eqref{val:recompile-goal}.

For the first assertion, the original auxiliary profile is fixed across
terminal histories.  Averaging its logarithmic charges and its pairing
volumes over those histories therefore gives
\[
 \sum_c w_c\log R_c=\log R,\qquad
 \sum_c w_c\log m_c=\log v.
\]
If the old rank bound charges unused resources or is larger than the product
rank, the first equality can instead be replaced by $\le$, which only
strengthens the conclusion below.  Non-GROUP operations have no additional
auxiliary charge or byproduct.  From \eqref{val:new-group-rates} and
\eqref{val:new-label-count}, the limiting improvement of the unnormalized
$\gamma$-score is at least
\begin{align*}
 &\sum_c w_c\bigl[
       \gamma\log R_c-\log m_c-(\gamma-1)h_c\bigr]\\
 &\quad=\sum_c w_c\bigl[
       (\log R_c-\log m_c)
       +(\gamma-1)(\log R_c-h_c)\bigr]\ge0,
\end{align*}
by \eqref{val:group-inequalities}.  Dividing by $n$ proves the first
assertion.  The finite common multiples needed above exist because the
expanded tree, its rational laws, and its old block sizes were fixed first.
\end{proof}

\begin{proposition}[Monotonicity]
\label{prop:monotonicity}
For $1\le\gamma_0\le\gamma_1$ and every $\mathsf W$ law $q$,
\[
 U^W_{\gamma_1}(q)\ge U^W_{\gamma_0}(q).
\]
\end{proposition}

\begin{proof}
Apply Lemma~\ref{val:recompile} to any finite recipe scored at $\gamma_0$.
On each of the recompiled realizations,
\[
 u_{\gamma_1}-u_{\gamma_0}
       =(\gamma_1-\gamma_0)(D-A).
\]
The second assertion of that lemma makes the limiting difference
nonnegative, while its first assertion preserves the old $\gamma_0$-score.
Thus the closure at $\gamma_1$ contains a value at least as large at the old
law.  Approximate an arbitrary pair in the $\gamma_0$ closed attainable set
by finite recipes and take a diagonal sequence of their recompilations.
Their laws converge to the same $q$; closedness proves the claim.
\end{proof}

\subsection{A direct matching lower bound}
\label{val:matching-section}

The first boundary estimate separates original scalar components without
an auxiliary tensor.  The matching argument uses progression-free hashing
and collision removal as in \citet[Section~6]{cw}; we give the details
with the bank normalization needed here.

\begin{lemma}[Progression-free labels]
\label{val:ap-set}
For all sufficiently large integers $L$ there is a set
$J\subseteq\{1,\ldots,L\}$ containing no nonconstant three-term arithmetic
progression and satisfying
\[
 \log|J|=\log L-O(\sqrt{\log L}).
\]
For primes $P\to\infty$, a subset of $\mathbb F_P$ of size
$P\exp(-O(\sqrt{\log P}))$ has the same property for progressions modulo
$P$.
\end{lemma}

\begin{proof}
This is the sphere construction of \citet{behrend}.  Put
$d=\lfloor\sqrt{\log L}\rfloor$ and
$m=\lfloor L^{1/d}/2\rfloor$.  Encode vectors in
$\{0,\ldots,m-1\}^d$ as integers in base $2m$.  Their encodings are
nonnegative and smaller than $(2m)^d\le L$.  There are at most
$d(m-1)^2+1$ possible squared norms, so one norm shell contains at least
$m^d/(d(m-1)^2+1)$ vectors.  Its logarithmic size is
$\log L-O(\sqrt{\log L})$: both the base overhead $d\log2$ and the
norm-shell loss $O(\log d+\log m)$ have that order.

If three encodings satisfy $a+c=2b$, digitwise addition on either side
produces no carries, so their vectors satisfy the same midpoint relation.
If all three have the same squared norm, writing the first and third as
$b\pm u$ gives $\|b+u\|^2+\|b-u\|^2=2\|b\|^2+2\|u\|^2$ and forces
$u=0$.  The three vectors coincide.  Translating all encodings by one puts
them in $\{1,\ldots,L\}$ without changing this property.

For a prime $P$, apply the integer construction inside
$\{1,\ldots,\lfloor P/4\rfloor\}$.  An integer of the form $a+c-2b$
from this interval has absolute value less than $P$.  If it vanishes modulo
$P$, it is zero over the integers.  This gives the asserted modular set.
\end{proof}

\begin{proposition}[Banked marginal matching]
\label{prop:matching}
Let $\alpha\in\Delta(\mathcal S)$, and let
$\alpha_X,\alpha_Y,\alpha_Z$ be its three canonical local-grade marginals.
For every $0<\beta\le3$,
\begin{equation}
 F_\beta(\alpha)\ge H_{\mathrm{marg}}(\alpha)
 :=\frac{H(\alpha_X)+H(\alpha_Y)+H(\alpha_Z)}3.
\label{val:matching-bound}
\end{equation}
The constructions proving this bound are strong realizations with no
auxiliary input and no additional matrix multiplication volume.
\end{proposition}

\begin{proof}
First take $\alpha$ rational and $N$ through its admissible denominators.
Write $p_i=\alpha(2e_i)$ and $E_i=\alpha(\mathbf1-e_i)$.  On canonical
side $i$, the frequency of grade two is exactly $p_i$.  If $b_i$ is the
frequency of grade one, then
\[
 b_i=E_{i+1}+E_{i+2},\qquad
 E_i=\tfrac12(b_{i+1}+b_{i+2}-b_i),
\]
with subscripts modulo three.  Thus the three marginal types determine the
whole joint type.  Local restrictions to those types, separately in each
orientation bank, retain precisely the original words of canonical law
$\alpha$ in every bank.  No nonlocal joint-type filter is used.

View the retained support as a three-partite hypergraph.  A vertex is a
complete retained local index word on one physical side; an edge is a
retained triple of local words.  There are
\[
 B_N=\left(\frac{N!}{\prod_{s\in\mathcal S}(N\alpha_s)!}\right)^6
\]
edges.  Each physical side sees each canonical marginal in exactly two
banks, so all three vertex parts have the same size
\[
 V_N=\prod_{i=0}^2
 \left(\frac{N!}{\prod_{a=0}^2(N\alpha_i(a))!}\right)^2.
\]
Consequently, with $n=6N$,
\begin{equation}
 \log B_N=nH(\alpha)+O(\log N),\qquad
 \log V_N=nH_{\mathrm{marg}}(\alpha)+O(\log N).
\label{val:matching-counts}
\end{equation}
Permuting positions within each bank acts transitively on each local
exact-type vertex set and preserves the edge set.  Every vertex therefore
has degree $\Delta_N=B_N/V_N$.  Two vertices on different sides determine
at most one edge, since the original grades sum to two coordinatewise.

Choose a prime $P>3$ between $X_N$ and $2X_N$, where
$X_N=\exp(\sqrt N)\max(3,\Delta_N)$; integer rounding or enlarging $X_N$
by a fixed amount is immaterial.  Such a prime exists for all sufficiently
large $N$ by Bertrand's postulate.  Let $J\subset\mathbb F_P$ be the set
of Lemma~\ref{val:ap-set}.  Choose independent uniform field elements
$d_X,d_Y$ and $w_1,\ldots,w_n$, and for the three physical index words
$I,J',K$ define local hashes
\begin{align*}
 h_X(I)&=d_X+\sum_t w_t I_t,\\
 h_Y(J')&=d_Y+\sum_t w_t J'_t,\\
 h_Z(K)&=\tfrac12\left(d_X+d_Y+\sum_t w_t(2-K_t)\right).
\end{align*}
The symbol $J'$ here is an index word; the progression-free set is $J$.
On every edge $h_X+h_Y=2h_Z$.  Retaining variables whose hashes lie in
$J$ therefore forces all three hashes on a surviving edge to coincide.
A fixed edge survives with probability $|J|/P^2$.

For clarity, condition on this edge surviving with common hash $a\in J$.
The equations fix $d_X$ and $d_Y$ once the vector $(w_t)$ is chosen; they
leave that vector uniformly distributed.  For a different edge sharing its
$X$ vertex, survival imposes the nonzero linear equation
$\sum_t w_t(J''_t-J'_t)=0$.  Its conditional probability is $1/P$.
The same reasoning holds for a shared $Y$ vertex.  For a shared $Z$ vertex,
tightness gives $I'+J''=I+J'$; survival is equivalent to
$\sum_t w_t(I'_t-I_t)=0$, again of probability $1/P$.
The difference vectors are nonzero because two shared vertices would force
the edges to agree.  Their entries belong to $\{-2,-1,0,1,2\}$, so they
remain nonzero modulo $P$.

There are fewer than $3\Delta_N$ neighboring edges.  A union bound shows
that the expected number of isolated surviving edges is at least
\begin{equation}
 \frac{B_N|J|}{P^2}
       \left(1-\frac{3\Delta_N}{P}\right)
       =V_N\exp(-o(N)).
\label{val:isolated-count}
\end{equation}
Here $\log(P/\Delta_N)=O(\sqrt N)$,
$\log(P/|J|)=O(\sqrt{\log P})=o(N)$, and
$3\Delta_N/P\to0$.  Thus some hash choice attains this many isolated edges.
Delete every vertex not belonging to an isolated surviving edge.  No extra
edge can remain: an extra edge touching any retained vertex would contradict
that vertex's edge being isolated in the filtered hypergraph.  The output
is therefore a direct sum of scalar original components.

The restriction preserves each chosen original component, with any
unresolved tensor attached to it, unchanged.  All chosen words have the
same exact canonical type in every bank.  Thus it is a strong realization
with $A=M=0$ and
$D\ge H_{\mathrm{marg}}(\alpha)-o(1)$ by
\eqref{val:matching-counts}--\eqref{val:isolated-count}.
Passing to the attainable closure proves \eqref{val:matching-bound} for
rational laws.  Approximation by rational laws, including on faces, and
continuity prove the general statement.
\end{proof}

\subsection{The five-component face}
\label{val:face-section}

On one face the six-state tensor has a genuine three-state $\mathsf W$ partition.
The remaining distinctions within its coarse groups must be separated only
after that partition has been resolved.  This gives the second boundary
estimate used in the certificate.

\begin{proposition}[$\mathsf W$ constructions on a face]
\label{prop:face}
Suppose $\alpha=(p_0,p_1,p_2,E_0,E_1,E_2)\in\Delta(\mathcal S)$ satisfies
$E_1=0$.  Put
\begin{align}
 Q&=(p_2+E_0,\ p_0+E_2,\ p_1),\label{val:face-Q}\\
 J(\alpha)&=\en(p_2)+\en(E_0)+\en(p_0)+\en(E_2)
                  -\en(Q_0)-\en(Q_1),\label{val:face-J}
\end{align}
where $\en(s)=-s\log s$ and $\en(0)=0$.  Then for $0<\beta\le3$,
\begin{equation}
 F_\beta(\alpha)\ge
 \frac\beta3\bigl(U^W_{3/\beta}(Q)+J(\alpha)\bigr).
\label{val:face-bound}
\end{equation}
In particular, if $1\le\gamma_0\le3/\beta$, the same lower bound holds
with $U^W_{\gamma_0}$ in place of $U^W_{3/\beta}$.
\end{proposition}

\begin{proof}
Scale $x_1,z_1,y_0,y_2$ by a degeneration parameter.  The five terms other
than 101 have degree one, whereas 101 has degree three.  Dividing by the
parameter and taking its leading coefficient leaves
\[
 A=x_0(y_0z_2+y_1z_1),\qquad
 B=(x_2y_0+x_1y_1)z_0,\qquad C=x_0y_2z_0.
\]
The local predicates $z>0$, $x>0$, and $y=2$ identify, respectively,
$A,B,C$, and exactly one holds on each retained term.  These are $\mathsf W$ flags,
with the canonical coarse law $Q$ in \eqref{val:face-Q}.  The fixed
permutation identifying the three active parties with the three coarse
labels is applied to the physical banks; it bijects their orientations.

Apply a strong $\mathsf W$ recipe to these coarse occurrences.  Its unresolved
fibers contain the internal binary choices in $A$ and $B$.  The strong
identity retains those choices without restricting them.  Once a coarse
word is direct, all sides know its $A$ and $B$ positions.  Within $A$, the
choice between 002 and 011 is independently readable on $Y$ and $Z$;
within $B$, the choice between 200 and 110 is independently readable on
$X$ and $Y$.  Apply exact conditional types and GROUP to the $A$ and $B$
pools separately in every original orientation bank.
Their total entropy per original occurrence is
\[
 Q_0 H\left(\frac{p_2}{Q_0},\frac{E_0}{Q_0}\right)
 +Q_1 H\left(\frac{p_0}{Q_1},\frac{E_2}{Q_1}\right)
 =J(\alpha),
\]
with an empty pool contributing zero.  Lemma~\ref{lem:group} gives this
additional logarithmic volume and matching direct multiplicity, with equal
auxiliary rank to leading order.  Exact types fix pool sizes across all
coarse histories, so the internal auxiliary product is charged once.
The final labels are distinct original scalar words of the prescribed
canonical bank law, and all their fibers are nonzero.

For a coarse realization scored at $\gamma=3/\beta$, its contribution to
the six-state score is $(\beta/3)u_\gamma$ by
\eqref{val:scores}.  The two internal pools contribute
$(\beta/3)J(\alpha)$ in the limit.  To obtain the supremum at an interior
face law, approximate its $\mathsf W$ value by finite recipes, correct their small
coarse-law discrepancies by Remark~\ref{val:dilution}, and choose rational
conditional types.  Every step is finite, with errors made arbitrarily
small before taking the attainable closure.  This proves
\eqref{val:face-bound} when the five present coordinates are positive.
Both sides are continuous on the closed face by
Proposition~\ref{prop:values} and continuity of conditional entropy, so
limits give the remaining cases.  The last assertion is
Proposition~\ref{prop:monotonicity}.
\end{proof}
